% M32 — Formal theorem statements for Conjecture M13.1 (paper-2 ANT-tier)
% f = 4.5.b.a (LMFDB CM newform of weight k=5, level N=4, character chi_4, CM by K=Q(i))
% Date: 2026-05-06. Hallu count: 85.

\section{Setup and notation}

Let $f \in S_5(\Gamma_0(4),\chi_4)$ denote the unique normalised cuspidal
newform of weight $k=5$ and level $N=4$ on the LMFDB label \texttt{4.5.b.a},
with nebentypus $\chi_4$ the unique non-trivial Dirichlet character modulo $4$.
The form $f$ has complex multiplication by the imaginary quadratic field
$K = \mathbb{Q}(i)$, with associated Hecke character $\psi$ of infinity type
$(k-1,0) = (4,0)$.

Let $p = 2$. Note $p$ \emph{ramifies} in $K$: $p\mathcal{O}_K = \mathfrak{p}^2$
where $\mathfrak{p} = (1+i)$. In particular $N = p^2 = 4$ and $f$ has Hecke
eigenvalue $a_p = a_2 = -4 = -p^{(k-1)/2}$, i.e.\ at the upper edge of the
bad-prime Ramanujan bound (the \emph{Steinberg-edge} of M13/M20).

Let $E:y^2 = x^3 + x$ be the elliptic curve of conductor $32$, $j=1728$,
with CM by $\mathbb{Z}[i]$. Its formal group $\hat{E}$ at $p=2$ has height
$2$ (supersingular reduction since $p=2$ ramifies in CM field).

Let $K_\infty^{\mathrm{anti}}/K$ be the anti-cyclotomic $\mathbb{Z}_p$-extension
of $K$ and $\Gamma = \mathrm{Gal}(K_\infty^{\mathrm{anti}}/K) \cong \mathbb{Z}_2$.
Fix a topological generator $\gamma \in \Gamma$. Let $\Lambda(\Gamma) :=
\mathbb{Z}_2[[\Gamma]]$ be the Iwasawa algebra and $\mathcal{D}(\Gamma,\mathbb{Z}_2)$
the space of $\mathbb{Z}_2$-valued locally-analytic distributions on $\Gamma$.

For $m \in \{1,2,3,4\}$ let $\chi_m$ be the unique Hecke character of $K$ of
infinity type $(m, k-1-m)$ and trivial finite order at $\mathfrak{p}$. The
Damerell rationals
\[
\alpha_1 = \tfrac{1}{10},\quad \alpha_2 = \tfrac{1}{12},\quad
\alpha_3 = \tfrac{1}{24},\quad \alpha_4 = \tfrac{1}{60}
\]
record the algebraic part of $L(f,\chi_m)/\Omega^{2m}$ at the
Chowla–Selberg period $\Omega = \Omega(K)$.

\section{Theorem M13.1.A — Existence of $L_2^\pm(f)$ via Kriz Hodge filtration}

\begin{conjecture}[M13.1.A]
\label{conj:M131A}
Let $f$, $K$, $p=2$, $\Gamma$, $\Omega_2$ be as in the Setup, where
$\Omega_2 \in \mathbb{Q}_2^\times$ is the $p$-adic period attached to a
generator of the Hodge-filtration line $\mathrm{Fil}^4 H^1_{\mathrm{dR}}(E)/\mathbb{Q}_2$
via the Maass–Shimura operator (cf.\ Kriz~2021, Ch.~3). Then there exist
two distributions
\[
L_2^+(f),\; L_2^-(f) \;\in\; \mathcal{D}(\Gamma,\mathbb{Z}_2)
\]
satisfying the interpolation property: for all $m \in \{1,2,3,4\}$,
\[
\boxed{\;
\int_\Gamma \chi_m\, dL_2^\pm(f) \;=\; E_2^\pm(f,m) \cdot
\frac{\alpha_m \pm \alpha_{k-m}}{2} \cdot \Omega_2^{2m}
\;}
\]
where $E_2^\pm(f,m) \in \mathbb{Q}_2^\times$ is the Euler-type compensation factor
\[
E_2^\pm(f,m) \;=\; (-p^{m-1})\,(1+p^{m-3})
\]
(the F1 formula of M22, derived after Frobenius zero-root degeneracy
compensation; \emph{[TBD: prove A3 — explicit Euler-factor derivation]}).
\end{conjecture}

\textbf{[TBD: prove]}: A1 (Kriz extension to ramified $p\mid N$), A2 (Fan–Wan
ramified principal series at $p=2$ for 4.5.b.a), A3 (explicit $E_2^\pm$ formula),
A4 ($\Omega_2$ from formal group of $E$ at $p=2$ ramified).

\section{Theorem M13.1.B — Boundedness via Iwasawa-log convention}

\begin{conjecture}[M13.1.B]
\label{conj:M131B}
Assume Conjecture~\ref{conj:M131A}. In the Iwasawa logarithm convention
$\log_p(p) := 0$, one has
\[
\log_p(\alpha) \;=\; \log_p(-4) \;=\; 0 \quad (\text{since } \log_p(-1) = \log_p(p) = 0).
\]
Define the renormaliser
$\omega(\gamma) := \log_p(\gamma)/(\gamma - 1) \in \mathrm{Frac}\,\Lambda(\Gamma)$.
Then the renormalised distributions
\[
\widetilde{L}_2^\pm(f) := \omega(\gamma) \cdot L_2^\pm(f)
\;\in\; \Lambda(\Gamma)\otimes_{\mathbb{Z}_2}\mathbb{Q}_2
\]
are \emph{bounded measures}; equivalently, $\widetilde{L}_2^\pm(f) \in
\Lambda(\Gamma)$ up to a fixed $\mathbb{Q}_2^\times$ scalar.
\end{conjecture}

\textbf{[TBD: prove]}: B5 (admissibility of $\log_p(-4)=0$ choice), B6
(rigorous bounded-measure proof analogous to Pollack 2003 Lemma~5.1).

\section{Theorem M13.1.C(i) — F1 strict monotone $v_2$-pattern (UNCONDITIONAL)}

\begin{theorem}[M13.1.C(i); sympy-verified M22]
\label{thm:M131Cmono}
Define the F1 renormalisation
\[
\alpha_m^{\mathrm{ren}} := \alpha_m \cdot (-p^{m-1}) \cdot (1 + p^{m-3}),
\quad p=2,\ m \in \{1,2,3,4\}.
\]
Then
\[
(\alpha_1^{\mathrm{ren}},\alpha_2^{\mathrm{ren}},\alpha_3^{\mathrm{ren}},
\alpha_4^{\mathrm{ren}}) = \bigl(-\tfrac{1}{8},\,-\tfrac{1}{4},\,-\tfrac{1}{3},\,-\tfrac{2}{5}\bigr),
\]
with strict monotone $2$-adic valuations
\[
\bigl(v_2(\alpha_m^{\mathrm{ren}})\bigr)_{m=1}^{4}
= (-3,\,-2,\,0,\,+1).
\]
\end{theorem}

\begin{proof}
Direct computation, sympy-verified in M22.
\end{proof}

\begin{conjecture}[M13.1.C(ii) — embedding into $L_2^\pm$]
\label{conj:M131Cembed}
Assume Conjecture~\ref{conj:M131A}. The Damerell rationals embed via
$\int_\Gamma \chi_m\,dL_2^\pm(f) / \Omega_2^{2m}$, so the integrality
property $v_2 \in \{-3,-2,0,+1\}$ corresponds to the $\mathbb{Z}_2$-valuation
of the distribution evaluated at $\chi_m$. \textbf{[TBD: prove C7, C8]}.
\end{conjecture}

\section{Theorem M13.1.D(i) — Pair-sum rationality requires square level (UNCONDITIONAL)}

\begin{theorem}[M13.1.D(i); classical functional equation]
\label{thm:M131Dsquare}
Let $g$ be a CM newform of weight $k$, level $N$, root number $\varepsilon$,
self-dual, with Damerell rationals $\alpha_m(g)$. The functional-equation
ratio
\[
R(m;k,N,\varepsilon) := \frac{\alpha_m(g)}{\alpha_{k-m}(g)} =
\varepsilon \cdot N^{(k-2m)/2} \cdot 2^{2m-k} \cdot
\frac{\Gamma(k-m)}{\Gamma(m)}
\]
is rational over $\mathbb{Q}$ if and only if $N^{(k-2m)/2} \in \mathbb{Q}$,
which (for odd $k$ and $m \neq k/2$) holds if and only if $N$ is a perfect
square. Consequently, the pair-sum $\alpha_m(g) + \alpha_{k-m}(g) \in \mathbb{Q}$
\emph{only if} $N$ is a perfect square.
\end{theorem}

\begin{proof}
Classical functional equation $\Lambda(s) = \varepsilon \Lambda(k-s)$
applied at $s=m$ vs $s=k-m$. Verified in 9-form LMFDB sample (M21):
4.5.b.a ($N=4$ square, pair-sum $1/8$), 36.5.d.a, 64.5.c.a, 100.5.b.a
(square but $\alpha_2 \neq 1/12$), 27.5.b.a, 81.5.d.a, 3.7.b.a, 27.3.b.a,
12.3.c.a (non-square or wrong CM, irrational pair-sums).
\end{proof}

\section{Theorem M13.1.D(ii) — Steinberg-edge $\Leftrightarrow$ $p^2\mid N$ (CONDITIONAL)}

\begin{theorem}[M13.1.D(ii); conditional on D10]
\label{thm:M131Dsteinberg}
Let $g$ be a CM newform with CM by imaginary quadratic field $K$ and Hecke
character $\psi$ unramified at $\mathfrak{p}\mid p$. Then $|a_p(g)| = p^{(k-1)/2}$
(Steinberg-edge condition) holds iff $p$ ramifies in $K$ and the Hecke-character
induction satisfies $a_p(g) = \psi(\mathfrak{p})$ with $|\psi(\mathfrak{p})|^2
= N(\mathfrak{p})^{k-1} = p^{k-1}$. This forces $p^2 \mid N$.
\textbf{[TBD: prove D10 unconditionally via local Langlands]}.
\end{theorem}

\section{Conjecture M13.1.D(iii) — Uniqueness of 4.5.b.a}

\begin{conjecture}[M13.1.D(iii)]
\label{conj:M131Dunique}
Among CM-by-$\mathbb{Q}(i)$ newforms of weight $5$ at perfect-square level
$N \leq 100$ with rational Damerell pair-sums summing to a pure $2$-power,
$f = $ \texttt{4.5.b.a} is the unique such form (verified in 9-form LMFDB
sample of M21). \textbf{[TBD: prove D11 — extend to all square levels]}.
\end{conjecture}

\section*{Companion remark (M27 / Hodge)}

The motive $M(f)$ of $f = $ \texttt{4.5.b.a} is a direct summand of $H^4$
of a power of the CM elliptic curve $E:y^2=x^3+x$, so the Hodge conjecture
for $M(f)$ is already a theorem via Pohlmann 1968 + Tate–Imai–Murty
correspondence (cf.\ M27 SUMMARY). A separate short companion note
(8--12pp, *Research in Number Theory*) addresses the Beilinson-regulator
analogue of Conjecture~\ref{conj:M131A} at $s=k=5$.
